\scanpage{148}
vu que
\[
g^{-1}=\begin{pmatrix}d&-b\\-c&a\end{pmatrix}
\]
la formule
\[
\tau\circ g^{-1}=\frac{a\tau-b}{-c\tau+d}. \tag*{(11)}
\]
Introduisons l'élément
\[
\left\{\begin{aligned}
&\tau_\infty=\begin{pmatrix}-1&0\\0&+1\end{pmatrix}\in\unsure{Gl|GL}(2,\mathbb Z)\qquad\text{d'où}\\
&\tau_\infty\begin{pmatrix}A&B\\C&D\end{pmatrix}\tau_\infty^{-1}
=\tau_\infty\left(\begin{pmatrix}A&B\\C&D\end{pmatrix}\right)
=\begin{pmatrix}A&-B\\-C&D\end{pmatrix}\\
&\tau_\infty\left(\begin{pmatrix}a&b\\c&d\end{pmatrix}^{-1}\right)
=\begin{pmatrix}d&b\\c&a\end{pmatrix}\qquad\substack{\text{``échange de}\\\text{$a$ et de $d$''}}
\end{aligned}\right.,\tag*{(12)}
\]
on trouve
\[
\begin{aligned}
\tau\circ\tau_\infty(g^{-1})&=\frac{a\tau+b}{c\tau+d}
&\text{où }\ g&=\begin{pmatrix}a&b\\c&d\end{pmatrix}\\
&&\tau_\infty(g^{-1})&=\begin{pmatrix}d&b\\c&a\end{pmatrix}
\end{aligned}\tag*{(13)}
\]
Par la suite, quand on fait opérer $SL(2,\mathbb Z)$ sur $\mathcal D^+$, ce sera tjrs (sauf mention expresse) par la formule (13) -- ce n'est donc pas l'action ``modulaire'' la plus évidente, mais celle consacrée par l'usage, qui correspond à la précédente \inserted{(à droite)}, modifiée via le passage
\[
\underset{\substack{\Vert\\\left(\begin{smallmatrix}a&b\\c&d\end{smallmatrix}\right)}}g\longmapsto\tau_\infty(g^{-1})=\begin{pmatrix}d&b\\c&a\end{pmatrix}.
\]
% archive sheet number: 147
